\documentclass[11pt]{article}
\usepackage[T1]{fontenc}
\usepackage{amsmath,amssymb,geometry}
\geometry{margin=1in}
\begin{document}
\begin{center}
{\Large Normalized spectral support need not fill an interval}\\[4pt]
Disproof of Conjecture 00000003702\\[4pt]
ziangni-sys --- October 8, 2026
\end{center}
\paragraph{Scope of the claim.}
Both language versions state that the limiting measure has support
the interval $[0,2]$. We disprove this literal equality. The usual
statement that normalized Laplacian spectra are \emph{contained}
in $[0,2]$ is valid and is not contradicted here. The source imposes
no connectedness hypothesis. Our graphs have unbounded vertex count.

\paragraph{Actual graphs and normalized operator.}
For $m\geq0$, let $G_m$ be the disjoint union of $m+1$ copies of
$K_2$. More explicitly its vertices are $(c,i)$ with
$c\in\{0,\ldots,m\}$ and $i\in\{0,1\}$; adjacency holds exactly
when the copy labels agree and the endpoints differ. Every vertex
has its opposite endpoint as its unique neighbor, so every degree
is one and $|V(G_m)|=2(m+1)\to\infty$.
Consequently the genuine normalized Laplacian is
\[
 L_m=I-D^{-1/2}A_mD^{-1/2}=I-A_m,
 \qquad (L_mf)(c,i)=f(c,i)-f(c,1-i).
\]
The random-walk normalization $I-D^{-1}A_m$ gives the same operator
because $D=I$.

\paragraph{Complete spectrum.}
On each edge the symmetric vector $(1,1)$ has eigenvalue zero,
and the antisymmetric vector $(1,-1)$ has eigenvalue two. They
form a basis of that edge's two-dimensional function space.
Putting these bases on the disjoint copies gives a complete
$2(m+1)$-element eigenbasis of the actual graph operator.
Its invertible coordinates on copy $c$ are
\[
 a_c=\frac{f(c,0)+f(c,1)}2,\qquad
 b_c=\frac{f(c,0)-f(c,1)}2,
 \qquad f(c,i)=a_c+(-1)^i b_c.
\]
Thus there are exactly $m+1$ eigenvalues zero and $m+1$ eigenvalues
two, including multiplicities.

\paragraph{Empirical measures and weak limit.}
The ordinary normalized empirical spectral probability measure is
\[
 \mu_m=\frac1{2(m+1)}\sum_{c=0}^{m}\bigl(\delta_0+\delta_2\bigr)
      =\frac12\delta_0+\frac12\delta_2=:\mu.
\]
This is a constant sequence of actual empirical measures, so it
converges weakly to $\mu$; equivalently every bounded continuous
test function has identical integral against every $\mu_m$ and $\mu$.

\paragraph{Failure of full support.}
The open neighborhood $U=(1/2,3/2)$ of the point $1$ has
$\mu(U)=0$, since it contains neither atom. By the standard
neighborhood definition of topological support, $1$ is absent
from $\operatorname{supp}\mu$, although $1\in[0,2]$.
Indeed the support is exactly $\{0,2\}$. Therefore the limiting
support is not the full interval asserted in the conjecture.
This false clause suffices for the disproof even though this
particular sequence does converge.

\paragraph{Formal certificate.}
Lean constructs the genuine simple graphs, certifies every degree,
defines the normalized operator from actual neighbor sums and
constructs the complete eigenbasis via invertible coordinates.
It verifies all eigenvectors and both multiplicity counts, forms
the empirical Dirac measures, proves their exact constant value
and their weak convergence in the probability-measure topology.
The standard neighborhood support definition is written explicitly
because this Mathlib pin has no packaged measure-support function.
It verifies the zero-mass open neighborhood and refutes support
equality with $[0,2]$. Final audits use only standard Lean and
Mathlib axioms.
\end{document}
